\documentclass[12pt,a4paper]{article}

% --- Packages ---
\usepackage{geometry}
\geometry{top=2.5cm, bottom=2.5cm, left=2.5cm, right=2.5cm}
\usepackage{graphicx}
\usepackage{amsmath}
\usepackage{amssymb}
\usepackage{float}
\usepackage[hidelinks]{hyperref}
\usepackage{xeCJK} % 用於支援封面中文顯示，請務必使用 XeLaTeX 編譯
\usepackage{listings}
\usepackage{xcolor}

% --- Font Settings (Optional: can be adjusted based on preference) ---
\setCJKmainfont{DFKai-SB} % 如果在 Overleaf 上沒有標楷體，可以改成 \setCJKmainfont{nofonts} 或是 bsmi 等內建字體

\graphicspath{{figure/}}

\begin{document}

% ==========================================
% Title Page (Based on the provided cover PNG)
% ==========================================
\begin{titlepage}
    \centering
    \vspace*{2cm}
    
    {\Huge \textbf{Digital Signal Processing} \par}
    \vspace{0.5cm}
    {\Huge \textbf{in VLSI Design} \par}
    
    \vspace{2.5cm}
    
    {\Huge \textbf{Homework 2} \par}
    \vspace{0.5cm}
    {\Huge \textbf{FIR Filter} \par}
    
    \vspace{4cm}
    
    {\Large
    \begin{tabular}{ll}
    Name : & 陳映宬 \\
    Department : & Electrical Engineering \\
    Student ID : & B11107027 \\
    Instructor : & 蔡佩芸
    \end{tabular}
    \par}
    
    \vfill
    
    {\Large March 31 , 2026 \par}
\end{titlepage}

\newpage
\tableofcontents
\newpage

% ==========================================
% I. Filter Design and Spectrum Analysis
% ==========================================
\section{Filter Design and Spectrum Analysis (Step 1)}

\subsection{Impulse Response, Frequency Response, and Coefficient Calculation}
\begin{figure}[H]
    \centering
    \includegraphics[width=0.8\textwidth]{step1_impulse_response.png}
    \caption{Impulse Response of the 25-tap sinc-based FIR filter.}
\end{figure}

\begin{figure}[H]
    \centering
    \includegraphics[width=0.8\textwidth]{step1_frequency_response.png}
    \caption{Magnitude and Phase Response of the FIR filter.}
\end{figure}

\newpage
\paragraph{Coefficient Calculation:}
The 25-tap coefficients are defined as $h[n] = \text{sinc}(n/6)$ for $-12 \le n \le 12$, then shifted by 12 samples to obtain a causal sequence.

The quantized 17-bit signed Verilog representation of the coefficients is:
\begin{lstlisting}[language=Verilog, basicstyle=\ttfamily\small, frame=single]
h_coef[ 0] = 17'sh1FFFF; //  -0.00000000
h_coef[ 1] = 17'sh1F4E3; //  -0.08681179
h_coef[ 2] = 17'sh1EAD4; //  -0.16539867
h_coef[ 3] = 17'sh1E4D6; //  -0.21220659
h_coef[ 4] = 17'sh1E589; //  -0.20674834
h_coef[ 5] = 17'sh1EE89; //  -0.13641852
h_coef[ 6] = 17'sh00000; //   0.00000000
h_coef[ 7] = 17'sh01872; //   0.19098593
h_coef[ 8] = 17'sh034ED; //   0.41349667
h_coef[ 9] = 17'sh0517C; //   0.63661977
h_coef[10] = 17'sh069DA; //   0.82699334
h_coef[11] = 17'sh07A3B; //   0.95492966
h_coef[12] = 17'sh08000; //   1.00000000
h_coef[13] = 17'sh07A3B; //   0.95492966
h_coef[14] = 17'sh069DA; //   0.82699334
h_coef[15] = 17'sh0517C; //   0.63661977
h_coef[16] = 17'sh034ED; //   0.41349667
h_coef[17] = 17'sh01872; //   0.19098593
h_coef[18] = 17'sh00000; //   0.00000000
h_coef[19] = 17'sh1EE89; //  -0.13641852
h_coef[20] = 17'sh1E589; //  -0.20674834
h_coef[21] = 17'sh1E4D6; //  -0.21220659
h_coef[22] = 17'sh1EAD4; //  -0.16539867
h_coef[23] = 17'sh1F4E3; //  -0.08681179
h_coef[24] = 17'sh1FFFF; //  -0.00000000
\end{lstlisting}

\subsection{Filter Type and Analysis}
\paragraph{Filter Type Analysis:} 
This design acts as a \textcolor[HTML]{CC0000}{\textbf{Low-pass filter}}, as the sinc function in time domain corresponds to a rectangular window in frequency domain, inherently passing low frequencies and suppressing high ones.

% ==========================================
% II. Time-Domain Signal Filtering Verification
% ==========================================
\section{Time-Domain Signal Filtering Verification (Step 2)}

\subsection{Time-Domain Input and Output Waveforms}
\begin{figure}[H]
    \centering
    \includegraphics[width=0.8\textwidth]{step2_time_domain.png}
    \caption{Waveforms of the 144 input samples $x[n]$ and the corresponding output $y[n]$.}
\end{figure}

\subsection{Filtering Effect and Frequency Component Analysis}
\paragraph{Filtering Effect:}
$x[n]$ contains two components: $\sin(-2\pi n/128)$ (normalized freq $\approx 0.0156$, passband) and $\cos(2\pi n/4)$ (normalized freq $= 0.5$, stopband). After filtering, the low-frequency sine is retained while the high-frequency cosine is suppressed, as confirmed by the output waveforms above.

% ==========================================
% III. Quantization Error Analysis and Optimal Word-Length Determination
% ==========================================
\section{Quantization Error Analysis and Optimal Word-Length Determination (Step 3)}

\subsection{Four-Stage Optimal Word-Length Scanning}
\begin{figure}[H]
    \centering
    \includegraphics[width=0.8\textwidth]{step3_wordlength_decision.png}
    \caption{RMSE vs. word-lengths scanning from 9 bits to 20 bits.}
\end{figure}

Based on the requirement, the RMSE threshold for my student ID (odd number) is $2^{-11}$. The final fractional word-lengths determined to meet this criteria are:
\begin{itemize}
    \item Input word-length ($b_{in}$): 13 bits
    \item Coefficient word-length ($b_{coef}$): 15 bits
    \item Multiplier output word-length ($b_{mult}$): 17 bits
    \item Adder output word-length ($b_{add}$): 17 bits
\end{itemize}

\begin{table}[H]
    \centering
    \caption{Verilog word-length summary for each node in the direct form FIR filter.}
    \begin{tabular}{llr}
        \hline
        \textbf{Node} & \textbf{Format} & \textbf{Total bits} \\
        \hline
        Input ($x$)                      & 1 Sign + 2 Int + 13 Frac & 16 bits \\
        Coefficient ($h$)                & 1 Sign + 1 Int + 15 Frac & 17 bits \\
        Multiplier output ($x \times h$) & 1 Sign + 2 Int + 17 Frac & 20 bits \\
        Adder output (acc)               & 1 Sign + 3 Int + 17 Frac & 21 bits \\
        \hline
    \end{tabular}
\end{table}

\subsection{Direct Form Filter Block Diagram and Truncation Mechanism}
% \begin{figure}[H]
%     \centering
%     \includegraphics[width=0.8\textwidth]{direct_form_diagram.png}
%     \caption{Block diagram of the Direct Form FIR filter with word-lengths.}
% \end{figure}

\paragraph{Truncation Mechanism:}
To control the word-length growth, truncation is applied after each arithmetic node. Specifically, the theoretical extended bits from multiplications (sum of input and coefficient word-lengths) are truncated to the selected multiplier word-length; similarly, additions are truncated to the selected adder word-length. These truncation points are explicitly indicated by dashed boxes in the Direct Form block diagram above.

\subsection{Floating-Point vs. Fixed-Point Comparison: Time Domain}
\begin{figure}[H]
    \centering
    \includegraphics[width=0.8\textwidth]{step3_time_domain_comparison.png}
    \caption{Time-domain comparison and error between fixed-point and floating-point models.}
\end{figure}

\subsection{Floating-Point vs. Fixed-Point Comparison: Frequency Domain}
\begin{figure}[H]
    \centering
    \includegraphics[width=0.8\textwidth]{step3_freq_domain_comparison.png}
    \caption{Frequency response and passband magnitude error between fixed-point and floating-point models.}
\end{figure}

% ==========================================
% IV. Hardware Implementation and Single-Stage Verification
% ==========================================
\section{Hardware Implementation and Single-Stage Verification (Step 4)}

\subsection{Behavioral Simulation and Latency Explanation}
\begin{figure}[H]
    \centering
    \includegraphics[width=\textwidth]{step4_behav_orig_first20.png}
    \includegraphics[width=\textwidth]{step4_behav_orig_first20b.png}
    \caption{Behavioral simulation timing diagram --- first 20 samples.}
\end{figure}
\begin{figure}[H]
    \centering
    \includegraphics[width=\textwidth]{step4_behav_orig_last20.png}
    \includegraphics[width=\textwidth]{step4_behav_orig_last20b.png}
    \caption{Behavioral simulation timing diagram --- last 20 samples.}
\end{figure}
\begin{figure}[H]
    \centering
    \includegraphics[width=0.8\textwidth]{step4_hw_vs_float.png}
    \caption{Error plot between hardware behavioral outputs and MATLAB floating-point results.}
\end{figure}

\paragraph{Latency Explanation:}
In this design, one level of D Flip-Flop is added to both the input port and the output port as required by the specification. Therefore, the overall system latency is \textcolor[HTML]{CC0000}{1 clock cycle}.

\subsection{Circuit Synthesis and Critical Path Analysis}

\paragraph{A) Theoretical Critical Path:}
The theoretical critical path of the direct form FIR filter consists of \textcolor[HTML]{CC0000}{1 multiplier} and \textcolor[HTML]{CC0000}{24 adders}, as each output sample requires one multiplication followed by 24 sequential additions.

\paragraph{B) Post-Synthesis Critical Path:}
After synthesis, the actual critical path contains \textcolor[HTML]{CC0000}{1 multiplier (DSP48E1)} and \textcolor[HTML]{CC0000}{6 adders (CARRY4)}, with total delay = \textcolor[HTML]{CC0000}{13.484 ns} under timing constraint of \textcolor[HTML]{CC0000}{14 ns}. Although the theoretical critical path has 24 adders (linear accumulation), the synthesizer restructures them into a balanced adder tree of depth $\lceil \log_2 25 \rceil = 5$ levels, implemented as CARRY4 carry chains on the FPGA. Therefore only 6 CARRY4 cells remain on the critical path, which is the expected behavior of FPGA synthesis optimization.

\begin{figure}[H]
    \centering
    \includegraphics[width=0.8\textwidth]{step4_timing_report.png}
    \caption{Post-synthesis timing report of the critical path (Original design).}
\end{figure}

\paragraph{Timing Constraint:}
The final clock period constraint is set to $T_{clk} = \textcolor[HTML]{CC0000}{\textbf{14}} \text{ ns}$ (\textcolor[HTML]{CC0000}{71.43 MHz}). The rationale behind this value is: by iteratively shrinking the constraint, the synthesis tool reports a Timing Violation when $T_{clk}$ is less than this value. Thus, this setting represents the maximum operating frequency limit achievable by this design.

\subsection{Post-Synthesis Simulation}
\begin{figure}[H]
    \centering
    \includegraphics[width=\textwidth]{step4_post_orig_first20.png}
    \includegraphics[width=\textwidth]{step4_post_orig_first20b.png}
    \caption{Post-synthesis simulation timing diagram --- first 20 samples.}
\end{figure}
\begin{figure}[H]
    \centering
    \includegraphics[width=\textwidth]{step4_post_orig_last20.png}
    \includegraphics[width=\textwidth]{step4_post_orig_last20b.png}
    \caption{Post-synthesis simulation timing diagram --- last 20 samples.}
\end{figure}
\begin{figure}[H]
    \centering
    \includegraphics[width=0.8\textwidth]{step4_hw_vs_float.png}
    \caption{Error plot between post-synthesis outputs and MATLAB floating-point results (Original design).}
\end{figure}

The error plots of the behavioral simulation, post-synthesis simulation, and the MATLAB fixed-point model are nearly identical (RMSE = 4.5324e-04, Max $|$error$|$ = 8.2243e-04 for all three cases). This confirms that the hardware correctly replicates the fixed-point arithmetic defined in the MATLAB model, and that synthesis optimization does not introduce additional numerical errors.


% ==========================================
% V. 2-Level Parallel Processing Hardware Acceleration
% ==========================================
\section{2-Level Parallel Processing Hardware Acceleration (Step 5)}

\subsection{Parallel Processing Principle}
The core concept of 2-level Parallel Processing is to simultaneously compute two output samples, $y[2k]$ and $y[2k+1]$, utilizing twice the hardware resources to double the throughput. Based on the FIR filtering formula expansion:
\begin{align*}
    y[2k] &= \sum_{i=0}^{L-1} h[i]\, x[2k-i] \\
    y[2k+1] &= \sum_{i=0}^{L-1} h[i]\, x[2k+1-i]
\end{align*}
Both calculations share the identical coefficients $h[i]$, but with input samples shifted by one time step. By unrolling these two sets of operations into independent filter circuits, they can be output simultaneously within the same clock cycle, thereby doubling the throughput without increasing the critical path delay.

\subsection{Accelerated Architecture Design and Critical Path}
\begin{figure}[H]
    \centering
    \includegraphics[width=0.8\textwidth]{Parallel_Architecture.png}
    \caption{Block diagram of the 2-level parallel processing architecture.}
\end{figure}
% \begin{figure}[H]
%     \centering
%     \includegraphics[width=0.8\textwidth]{parallel_timing_report.png}
%     \caption{Synthesis report of the max delay for the accelerated design.}
% \end{figure}

The parallel design instantiates two independent direct-form FIR filter circuits computing $y[2k]$ and $y[2k+1]$ simultaneously. Since the two paths are independent, the critical path remains identical to the original design: \textcolor[HTML]{CC0000}{1 multiplier (DSP48E1) and 6 adders (CARRY4)}, with total delay = \textcolor[HTML]{CC0000}{13.484 ns}.

\begin{figure}[H]
    \centering
    \includegraphics[width=0.8\textwidth]{step5_timing_report.png}
    \caption{Post-synthesis timing report of the critical path (Parallel design).}
\end{figure}

The result is identical to the original design, confirming that parallelism does not increase the critical path delay.

\subsection{Accelerated Behavioral Simulation and Throughput Analysis}
\begin{figure}[H]
    \centering
    \includegraphics[width=\textwidth]{step5_behav_para_first20.png}
    \includegraphics[width=\textwidth]{step5_behav_para_first20b.png}
    \caption{Accelerated behavioral simulation timing diagram --- first 20 samples.}
\end{figure}
\begin{figure}[H]
    \centering
    \includegraphics[width=\textwidth]{step5_behav_para_last20.png}
    \includegraphics[width=\textwidth]{step5_behav_para_last20b.png}
    \caption{Accelerated behavioral simulation timing diagram --- last 20 samples.}
\end{figure}

The parallel design produces \textcolor[HTML]{CC0000}{2 outputs per clock cycle}, completing $y[0]$ to $y[143]$ within \textcolor[HTML]{CC0000}{$T_{out}/2 = 72 \times T_{clk}$}, compared to \textcolor[HTML]{CC0000}{$144 \times T_{clk}$} for the original design. The critical path delay remains unchanged since the two filter circuits operate independently in parallel.

\subsection{Accelerated Post-Synthesis Simulation}
\begin{figure}[H]
    \centering
    \includegraphics[width=\textwidth]{step5_post_para_first20.png}
    \includegraphics[width=\textwidth]{step5_post_para_first20b.png}
    \caption{Accelerated post-synthesis simulation timing diagram --- first 20 samples.}
\end{figure}
\begin{figure}[H]
    \centering
    \includegraphics[width=\textwidth]{step5_post_para_last20.png}
    \includegraphics[width=\textwidth]{step5_post_para_last20b.png}
    \caption{Accelerated post-synthesis simulation timing diagram --- last 20 samples.}
\end{figure}
\begin{figure}[H]
    \centering
    \includegraphics[width=0.8\textwidth]{step4_parallel_hw_vs_float.png}
    \caption{Error plot between post-synthesis outputs and MATLAB floating-point results (Parallel design).}
\end{figure}

\subsection{Hardware Cost (Area) Comparison Analysis}
\begin{table}[H]
    \centering
    \begin{tabular}{|l|c|c|c|}
    \hline
    \textbf{Resource} & \textbf{Original (Step 4)} & \textbf{Parallel (Step 5)} & \textbf{Ratio (Parallel/Original)} \\ \hline
    LUT  & 265 & 529 & 2.00$\times$ \\ \hline
    FF   & 423 & 460 & 1.09$\times$ \\ \hline
    DSP  & 20  & 40  & 2.00$\times$ \\ \hline
    IO   & 41  & 78  & 1.90$\times$ \\ \hline
    \end{tabular}
    \caption{FPGA resource utilization comparison.}
\end{table}

LUT and DSP doubled (2.00$\times$) as expected from duplicating the filter logic. FF increased only 1.09$\times$ since the shift register chain is largely shared. IO increased 1.90$\times$ due to the doubled input/output ports.


% % ==========================================
% % VI. Appendices and Submission Checklist
% % ==========================================
% \section{Appendices and Submission Checklist (Step 9)}
% The following files have been checked and uploaded to NTUCool:
% \begin{itemize}
%     \item[\checkmark] This PDF report containing all required results.
%     \item[\checkmark] Verilog RTL codes and Testbenches for Step 4 (Single-stage).
%     \item[\checkmark] Verilog RTL codes and Testbenches for Step 5 (Parallel processing).
%     \item[\checkmark] MATLAB/Python scripts for Step 1 to Step 3.
% \end{itemize}

\end{document}